% part: intuitionistic-logic
% chapter: introduction
% section: axiomatic-derivations

\documentclass[../../../include/open-logic-section]{subfiles}

\begin{document}

\olfileid{int}{int}{axd}

\olsection{स्वयंसिद्धकीय \usetoken{P}{derivation}}

अंतःप्रज्ञावादी विधानीय तर्कशास्त्रातील स्वयंसिद्धकीय
!!{derivation}s या संकल्पनात्मकदृष्ट्या सर्वांत सोप्या
आणि ऐतिहासिकदृष्ट्या पहिल्या !!{derivation}
पद्धती आहेत. त्यांचे कार्य अभिजात विधानीय
तर्कशास्त्रातल्याप्रमाणेच चालते.

\begin{defn}[!!^{derivability}]
जर $\Gamma$ हा $\Lang L$ मधील !!{formula}s चा संच
असेल, तर $\Gamma$ पासूनची \emph{!!{derivation}} म्हणजे
!!{formula}s चा सांत अनुक्रम $!A_1$, \dots,~$!A_n$
असा की प्रत्येक $i \le n$ साठी पुढीलपैकी एक अट
पूर्ण होते:
\begin{enumerate}
\item $!A_i \in \Gamma$; किंवा
\item $!A_i$ हे स्वयंसिद्धक आहे; किंवा
\item एखाद्या $!A_j$ आणि $!A_k$ पासून, जिथे $j < i$
  आणि $k < i$, विध्यात्मक अनुमानाने $!A_i$ निष्पन्न होते;
  म्हणजे $!A_k \ident !A_j \lif !A_i$.
\end{enumerate}
\end{defn}

\begin{defn}[स्वयंसिद्धके]
अंतःप्रज्ञावादी विधानीय तर्कशास्त्रातील
\emph{स्वयंसिद्धकांचा} संच $\PAx$ हा पुढील रूपांच्या
सर्व !!{formula}s चा बनलेला आहे:
\begin{align}
  & (!A \land !B) \lif !A \ollabel{ax:land1}\\
  & (!A \land !B) \lif !B \ollabel{ax:land2}\\
  & !A \lif (!B \lif (!A \land !B)) \ollabel{ax:land3}\\
  & !A \lif (!A \lor !B) \ollabel{ax:lor1}\\
  & !A \lif (!B \lor !A) \ollabel{ax:lor2}\\
  & (!A \lif !C) \lif ((!B \lif !C) \lif ((!A \lor !B) \lif !C)) \ollabel{ax:lor3}\\
  & !A \lif (!B \lif !A) \ollabel{ax:lif1}\\
  & (!A \lif (!B \lif !C)) \lif ((!A \lif !B) \lif (!A \lif !C)) \ollabel{ax:lif2}\\
  & \lfalse \lif !A \ollabel{ax:lfalse1}
\end{align}
\end{defn}

\begin{defn}[!!^{derivability}]
!!^a{formula}~$!A$ हे $\Gamma$ पासून \emph{!!{derivable}}
असते, म्हणजे $\Gamma \Proves !A$, जर $\Gamma$
पासून सुरू होऊन $!A$ वर संपणारी !!a{derivation} असेल.
\end{defn}

\begin{defn}[प्रमेये]
रिक्त संचापासून $!A$ ची !!a{derivation} असेल,
तर !!^a{formula}~$!A$ हे \emph{प्रमेय} असते.
$!A$ हे प्रमेय असेल तर आपण $\Proves !A$ लिहितो
आणि ते प्रमेय नसेल तर $\Proves/ !A$ लिहितो.
\end{defn}

\begin{prop}
  अंतःप्रज्ञावादी तर्कशास्त्राच्या स्वयंसिद्धकीय
  !!{derivation} पद्धतीत $\Gamma \Proves !A$
  असेल, तर अभिजात तर्कशास्त्रातही
  $\Gamma \Proves !A$ असते. विशेषतः, $!A$ हे
  अंतःप्रज्ञावादी प्रमेय असेल, तर ते अभिजात प्रमेयही असते.
\end{prop}

\begin{proof}
  प्रत्येक अंतःप्रज्ञावादी स्वयंसिद्धक हे अभिजात
  स्वयंसिद्धकही आहे. म्हणून अंतःप्रज्ञावादी
  तर्कशास्त्रातील प्रत्येक !!{derivation} ही अभिजात
  तर्कशास्त्रातीलही !!a{derivation} आहे.
\end{proof}

\end{document}
